\documentclass[12pt]{article}

%%\usepackage[english]{babel}
\usepackage{amsmath}
\usepackage{amssymb}
\usepackage{url}


\begin{document}

\title{Econometrics Notes IV: Elementary Statistics - Part 1: Results on Variance and Covariance}
\author{Artur Wegrzyn}
\maketitle

\abstract{This note presents derivations (or motivations, very heuristic at places) of some basic statistical facts that are invoked elsewhere in the notes. }

\tableofcontents



\section{Notation}
Sample $X_1, X_2, \dots, X_n$ is drawn independently from a distribution that
has finite mean and variance, i.e.:
\begin{eqnarray*}
X_i &\sim & i.i.d. \\
E(X_i) &=& E(X) = \mu = \mu_X < \infty \\
Var(X_i) &=& Var(X) = \sigma^2 = \sigma^2_X < \infty
\end{eqnarray*}

Where two-dimensional random variables are discussed, sample 
$\{ (X_1 ,Y_1), (X_2, Y_2), ..., (X_n, Y_n) \}$ is assumed to be drawn independently 
and the following assumptions are made:

\begin{eqnarray*}
(X_i, Y_i) &\sim & i.i.d. \\
E(X_i) &=& \mu_X < \infty \\
E(Y_i) &=& \mu_Y < \infty \\
Var(X_i) &=& = \sigma^2_X < \infty \\
Var(Y_i) &=& = \sigma^2_Y < \infty \\ 
Cov(X_i, Y_j) &=& \begin{cases}
\sigma_{XY} < \infty & \text{if } i=j \\ 
0 & \text{if } i \neq j
\end{cases}
\end{eqnarray*}



\section{Biased Sample Variance Rewritten}
Given $S^2$ defined as:
\begin{equation}
\label{biased_variance_estimator}
S^2 = \frac{1}{n}\sum_{i=1}^{n} (X_i - \overline{X})^2
\end{equation}
Note that:
\begin{eqnarray*}
\frac{1}{n}\sum_{i=1}^{n} (X_i - \overline{X})^2 &=& 
\frac{1}{n}\sum_{i=1}^{n}\left( 
(X_i - \mu) - (\overline{X} - \mu)
\right)^2 \\
&=& \frac{1}{n} \sum_{i=1}^{n}(X_i - \mu)^2 + 
\frac{1}{n} \sum_{i=1}^{n}(\overline{X} - \mu)^2 \\
&& - 2 (\overline{X} - \mu) \frac{1}{n} \sum_{i=1}^{n} (X_i - \mu) \\
&=& \frac{1}{n} \sum_{i=1}^{n}(X_i - \mu) + (\overline{X} - \mu)^2 
- 2 (\overline{X} - \mu)^2 + \\ && (\overline{X} - \mu)^2 \\
&=& \frac{1}{n} \sum_{i=1}^{n} (X_i - \mu)^2 - (\overline{X} - \mu)^2
\end{eqnarray*}
I.e.:
\begin{equation}
\label{biased_variance_estimator_decomposition}
\frac{1}{n} \sum_{i=1}^{n} (X_i - \overline{X})^2 = 
\frac{1}{n} \sum_{i=1}^{n} (X_i - \mu)^2 - (\overline{X} - \mu)^2
\end{equation}


\section{Biased Sample Covariance Rewritten}
\label{sec:biased_sample_covariance_rewritten}
Given $S_{XY}$ defined as:
\begin{equation}
\label{biased_variance_estimator}
S_{XY} = \frac{1}{n}\sum_{i=1}^{n} (X_i - \overline{X}) (Y_i - \overline{Y})
\end{equation}
Nothing special, just rewriting:
\begin{eqnarray*}
\frac{1}{n}\sum_{i=1}^{n} (X_i - \overline{X}) (Y_i - \overline{Y}) &=&
\frac{1}{n}\sum_{i=1}^{n}\left(
X_i Y_i - X_i \overline{Y} - \overline{X} Y_i + \overline{X} \cdot \overline{Y} 
\right) \\
&=& \frac{1}{n} \sum_{i=1}^{n} X_i Y_i - \overline{X} \cdot \overline{Y}
- \overline{X} \cdot \overline{Y} + \overline{X} \cdot \overline{Y} \\
&=& \frac{1}{n} \sum_{i=1}^{n} X_i Y_i - \overline{X} \cdot \overline{Y}
\end{eqnarray*}

Hence:
\begin{equation}
\frac{1}{n} \sum_{i=1}^{n} (X_i - \overline{X}) (Y_i - \overline{Y}) = 
\frac{1}{n} \sum_{i=1}^{n} X_i Y_i - \overline{X} \cdot \overline{Y}
\end{equation}

Which is sample analogon of well-known relation for expectations:
\begin{equation*}
Cov(X, Y) = E(X-E(X))(Y-E(Y)) = E(XY) - E(X) E(Y)
\end{equation*}


\section{Biased and Unbiased Variance Estimator}
Consider the expected value of the $S^2$ - expectation is applied to 
\eqref{biased_variance_estimator_decomposition}:
\begin{equation*}
E(S^2) = Var(X) - \frac{1}{n} Var(X) = \frac{n-1}{n} Var(X)
\end{equation*}
Hence: $S^2$ is a biased estimator of variance.\\
\indent Unbiased estimator of the variance $\tilde{S}^2$ is therefore given by:
\begin{equation*}
\tilde{S}^2 = \frac{n}{n-1} S^2 = 
\frac{1}{n-1} \sum_{i=1}^{n} (X_i - \overline{X})^2
\end{equation*}

\section{Biased and Unbiased Covariance Estimator}
Consider, in turn, the sample covariance estiamtor:
\begin{equation*}
S_{XY} = \frac{1}{n} \sum_{i=1}^{n}X_i Y_i - \overline{X} \cdot \overline{Y}
\end{equation*}

Taking expectation:
\begin{equation*}
E(S_{XY}) = \frac{1}{n} \sum_{i=1}^{n} E(X_i Y_i) - 
E(\overline{X} \cdot \overline{Y}) 
\end{equation*}

Rewriting minuend:
\begin{eqnarray*}
\frac{1}{n} \sum_{i=1}^{n} E(X_i Y_i) = 
\frac{1}{n} \sum_{i=1}^{n} \left(
\sigma_{XY} + \mu_X \mu_Y
\right) = \sigma_{XY} + \mu_X \mu_Y
\end{eqnarray*}


Rewriting subtrahend:
\begin{eqnarray*}
E(\overline{X} \cdot \overline{Y})  &=& 
\frac{1}{n^2}E \left(
\sum_{i=j} X_i Y_j + \sum_{i \neq j} X_i Y_j
\right) \\
&=&
\frac{1}{n^2} \sum_{i=1}^{n}E(X_i Y_i) +
\frac{1}{n^2} \sum_{i \neq j}E(X_i Y_j) \\
&=&
\frac{1}{n^2} \sum_{i=1}^{n} \left( \sigma_{XY} + \mu_X \mu_Y \right) +
\frac{1}{n^2} \sum_{i \neq j} \mu_X \mu_Y \\
&=&
\frac{1}{n} \left( \sigma_{XY} + \mu_X \mu_Y \right) +
\frac{n(n-1)}{n^2} \mu_X \mu_Y \\
&=&
\frac{1}{n}\sigma_{XY} + \frac{1}{n}\mu_X \mu_Y +
\frac{n-1}{n} \mu_X \mu_Y \\
&=&
\frac{1}{n}\sigma_{XY} + \mu_X \mu_Y 
\end{eqnarray*}

Plugging back the minuend and subtrahend into the formula for $E(S_{XY})$:
\begin{eqnarray*}
E(S_{XY}) = \sigma_{XY} + \mu_X \mu_Y - \frac{1}{n}\sigma_{XY} - \mu_X \mu_Y
\end{eqnarray*}

Briefly:
\begin{equation*}
E(S_{XY}) = \frac{n-1}{n}\sigma_{XY}
\end{equation*}

So that:
\begin{equation*}
\sigma_{XY} = \frac{n}{n-1} E(S_{XY}) = E(\frac{n}{n-1} S_{XY})
\end{equation*}

By section \ref{sec:biased_sample_covariance_rewritten} the biased covariance estimator
can be rewritten as:
\begin{equation*}
S_{XY} = \frac{1}{n} \sum_{i=1}^{n}(X_i - \overline{X})(Y_i - \overline{Y})
\end{equation*}

So the unbiased estimator of covariance is:
\begin{equation*}
\tilde{S}_{XY} = \frac{1}{n-1} \sum_{i=1}^{n} (X_i - \overline{X}) (Y_i - \overline{Y})
\end{equation*}



\section{Change of Variables for PDF of Two-Dimensional Variables}
In this section, a heuristic motivation of a formula for PDF of a two-dimensional random variable $(U, V)$ which is obtained from two-dimensional random variable 
$(X, Y)$ by means of a sufficiently regular transormation. 
It cannot be stressed enough, that this is \textbf{not} a strict proof of the result - rather: a motivation thereof under generous assumptions (which are for the most part sufficent for the purpose of these notes).

The CDF of a two-dimensional, continuous random variable $(X, Y)$ is:
\begin{equation*}
F_{X, Y}(x, y) = P(X \leq x, Y \leq y) = 
\int_{-\infty}^{y} \int_{-\infty}^{x} f_{X, Y}(s, t) ds dt 
\end{equation*}

Analogously, for $(U, V)$:
\begin{equation*}
F_{U, V}(u, v) = P(U \leq u, V \leq v) = 
\int_{-\infty}^{v} \int_{-\infty}^{u} f_{U, V}(s, t) ds dt
\end{equation*}


Suppose that $(X, Y)$ are related to $(U, V)$ by known transformation:
\begin{equation*}
\left\{
\begin{aligned}
x &= x(u, v) \\ 
y &= y(u, v)
\end{aligned}
\right.
\end{equation*}


Assume, further, that the transformation is invertible so the following relations can be obtained:
\begin{equation*}
\left\{
\begin{aligned}
u &= u(x, y) \\
v &= v(x, y)
\end{aligned}
\right.
\end{equation*}

Now, consider the region of $D$ of $Oxy$ plane such that:
\begin{equation*}
\left\{
\begin{aligned}
u(x, y) \leq u \\
v(x, y) \leq v
\end{aligned}
\right.
\end{equation*} 



Assuming sufficient regulatrity conditions one has:

\begin{equation*}
\left\{
\begin{aligned}
du = \frac{\partial u}{\partial x}dx + \frac{\partial u}{\partial y} dy \\
dv = \frac{\partial v}{\partial x}dx + \frac{\partial v}{\partial y} dy
\end{aligned}
\right.
\end{equation*}

Jacobian determinant of the transformation is: 

\begin{equation*}
\frac{D(u,v)}{D(x, y)} = 
\begin{vmatrix}
\frac{\partial u}{\partial x} & \frac{\partial u}{\partial y} \\ 
\frac{\partial v}{\partial x} & \frac{\partial v}{\partial y}
\end{vmatrix}
\end{equation*}

Now, consider probability of $(X, Y)$ falling within region defined by inequalities
$x < X \leq x + dx$ and  $d < Y \leq y + dy$:

\begin{equation*}
P(x < X \leq x + dx, y < Y \leq y + dy) = f_{X, Y}(x, y) dxdy
\end{equation*}

Since $(X, Y)$ and $(U, V)$ are related by the transformation mentioned above, 
let $dS$ be the the region on the $(U, V)$ plane obtained
due to infinitesimal increase of $X, Y$ by respectively $dx, dy$:

\begin{equation*}
dS = \left\{
(u,v): 
u(x, y) < u \leq u(x + dx, y + dy), 
v(x, y) < v \leq v(x + dx, y + dy)
\right\}
\end{equation*}


Then, letting $\vec{a}$ stand for the vector representing shift on the $(U, V)$
plane due to change of $X$ by $dx$ and 
letting $\vec{b}$ stand for the vector representing shift on the $(U, V)$
plane due to change of $Y$ by $dy$ we can write:

\begin{equation*}
\left\{
\begin{aligned}
\vec{a} = \left[
\frac{\partial u}{\partial x} dx, 
\frac{\partial v}{\partial x} dx
\right]^{T} \\ 
\vec{b} = \left[
\frac{\partial u}{\partial y} dy, 
\frac{\partial v}{\partial y} dy
\right]^{T}
\end{aligned}
\right.
\end{equation*}

Since on a 2-dimensional surface the area of a parallelogram spanned by 
two vectors is given by the absolute value of determinant of matrix formed from
components of these two rows as columns or rows, we obtain:
\begin{equation*}
|dS| = \begin{vmatrix}
\frac{\partial u}{\partial x} dx & \frac{\partial v}{\partial x} dx \\
\frac{\partial u}{\partial y} dy & \frac{\partial v}{\partial y} dy
\end{vmatrix} = 
\begin{vmatrix}
\frac{\partial u}{\partial x} & \frac{\partial v}{\partial x} \\
\frac{\partial u}{\partial y} & \frac{\partial v}{\partial y}
\end{vmatrix} dx dy =
\frac{D(u, v)}{D(x, y)}
\end{equation*}
Now, consider the following probability of the following event:
\begin{equation*}
P((U, V) \in dS) = 
f_{U, V}(u, v) |dS| = f_{U, V}(u, v) 
\frac{D(u, v)}{D(x, y)}dxdy
\end{equation*}

At the same time, by definition of $dS$ we also have:
\begin{equation*}
P((U, V) \in dS) = P(x < X \leq x + dx, y < Y \leq y + dy) = f_{X, Y}(x, y) dxdy
\end{equation*}

Comparison of RHS-es of last two equations yields:
\begin{equation*}
f_{X,Y}(x,y) = f_{U, V}(u,v) \frac{D(u, v)}{D(x, y)}
\end{equation*}
Given the invertibility of the transformation, this can be rewritten for $f_{U, V}$ as:
\begin{equation*}
f_{U, V}(u, v) = f_{X, Y}(x(u, v), y(u, v)) \frac{D(x, y)}{D(u, v)}
\end{equation*}

And the CDF of $(U, V)$ is:
\begin{eqnarray*}
F_{U, V}(u_0, v_0) &=& P(U \leq u_0, V \leq v_0) \\ 
&=& \int_{-\infty}^{v_0} \int_{-\infty}^{u_0} f_{U, V}(u, v) du dv \\
&=& \int_{-\infty}^{v_0} \int_{-\infty}^{u_0} f_{X, Y}(x(u, v), y(u, v)) 
\frac{D(x, y)}{D(u, v)} du dv
\end{eqnarray*}

\section{$\chi^2$ Distribution}
In this and subsequent section we consider a sample $X_1, X_2, \dots, X_n$ 
that is drawn independently from a normal 
distribution with mean $\mu$ and variance $\sigma^2$:
\begin{equation*}
X_i \overset{\text{i.i.d.}}{\sim} N(\mu, \sigma^2)
\end{equation*}


\subsection{Fisher's Theorem}
In this and subsequent section use is made of the Fisher's theorem that states 
that given the sample discussed, the sample mean and sample variance are independent 
random variables. I.e. given:
\begin{equation*}
\overline{X} = \frac{1}{n} \sum_{i=1}^{n} X_i
\end{equation*}

and:
\begin{equation*}
\tilde{S}^2 = \frac{1}{n-1} \sum_{i=1}^{n} (X_i - \overline{X})^2
\end{equation*}

we have:
\begin{equation*}
\overline{X} \perp\!\!\!\perp \tilde{S}^2
\end{equation*}

Proof of this theorem will be given in a separate note. For now, note that the symbol 
$\perp\!\!\!\perp$ is commonly used to denote independence of random variables. Also, 
recall that in case of normal distriubtion independence and orthogonality 
are equivalent - this will be a topic of yet separate note.


\section{Gamma Distribution}
Let $Y_1, Y_2, \dots, Y_n$ be a sample drawn indepentently from an exponential 
distribution with expected value $1 / \lambda$:

\section{Square of Standard Normal Distribution}
Suppose that $U$ follows standard normal distribution:
\begin{equation*}
U \sim N(0, 1)
\end{equation*}

Consider $Z = U^2$. Then:
\begin{equation*}
P(Z \leq z) = P(U^2 \leq z) = P(-\sqrt{z} < U \leq \sqrt{z})
\end{equation*}

Hence:
\begin{equation*}
f_{Z}(z) = \frac{dF_{Z}(z)}{dz} = \frac{d}{dz} P(-\sqrt{z} < U \leq \sqrt{z})
\end{equation*}

A little bit of algebra yields:
\begin{equation*}
f_{Z}(z) = 
\end{equation*}



\subsection{Statistical Motivation}
Given the sample, consider the biased variance estimator:
\begin{equation*}
S_{XY} = \frac{1}{n} \sum_{i=1}^{n} (X_i - \overline{X})^2 
\end{equation*}

The aim is to find its distribution. As previously proved, this estimator can be rewritten as:
\begin{equation*}
S_{XY} = \frac{1}{n} \sum_{i=1}^{n} (X_i - \mu)^2 - (\overline{X} - \mu)^2
\end{equation*}


\subsection{Derivation of PDF}

\section{Derivation of t-Student Distribution PDF}
\subsection{t-Student Distribution - Statistical Motivation}

\subsection{Derivation of the PDF}

\section{Derivation of F-Snedecor Distribution PDF}
\subsection{F-Snedecor Distribution - Statistical Motivation}
\subsection{Derivation of the PDF}

\bibliography{tsa_bibliography}
\bibliographystyle{ieeetr}

\end{document}
